%% ==========================================================================
%%  SECTION 3 --- CRYPTOGRAPHIC BUILDING BLOCKS                 [Track T1]
%%
%%  Job: give the reader exactly the primitives needed to follow \S4 and \S5,
%%  and land ONE argument -- that function secret sharing is what makes both
%%  halves of the problem tractable, with the SAME construction serving as a
%%  comparison gate in training and as a read layer in retrieval.
%%  ~2 pages. Resist the urge to write a crypto textbook.
%% ==========================================================================
\section{Cryptographic Building Blocks}
\label{sec:prim}

\todo{Frame the section by its punchline, stated up front: one primitive
      (FSS) turns out to be load-bearing on both sides of the problem. Every
      subsection below is evidence for that claim.}

\subsection{Secret sharing}
\label{sec:prim:sharing}

\todo{Additive sharing over a ring; then 2-out-of-3 REPLICATED sharing
      (Araki et al., CCS'16), which is what \nudgeplain{} uses: $x = x_0+x_1+x_2$
      with $\party{i}$ holding $(x_i, x_{i+1})$. Introduce $\shr{x}$ notation
      (macro in preamble.tex).
      Then the ``3PC with PRF'' model --- each PAIR of parties shares a PRF
      key, so correlated randomness costs no communication.}

\begin{observation}
\label{obs:prim:matvec-free}
\todo{THE key fact of the whole survey, stated as a numbered observation so
      later sections can point at it. Under replicated sharing, a degree-two
      function costs one round and $3\dime\bits$ bits, and consequently a
      shared-matrix $\times$ shared-vector product communicates in proportion
      to the largest intermediate VECTOR, not to the size of the input
      MATRIX. \nudgeplain{} Thm 4.2. \verify{quote the exact bound}}
\end{observation}

\subsection{Function secret sharing and distributed point functions}
\label{sec:prim:fss}

\todo{Boyle--Gilboa--Ishai (EUROCRYPT'15, CCS'16). Define a $(2,2)$-DPF:
      $\Gen(\alpha,\beta) \to (k_0,k_1)$ with
      $\Eval(k_0,x)-\Eval(k_1,x) = \beta$ iff $x=\alpha$, else $0$, and
      each key alone pseudorandom.
      The GGM-tree construction in a paragraph plus one figure. Key size
      $\dime(\seclen+2)+\bits$ bits --- for $\items=4096$ that is $\approx$260
      bytes against 32\,KB for a one-hot vector. THAT RATIO IS THE POINT; give
      it as a concrete number, not asymptotically.
      Cite the July 2026 DPF/FSS survey (arXiv 2607.27696) as the entry point,
      and mention the current frontier: shorter information-theoretic keys
      (2604.24385), multi-party DDH constructions (2603.17453), GPU FSS
      (FuseFSS 2606.09551).}

\begin{figure}[t]
\centering
\todo{Figure: the GGM tree, showing the two parties' paths diverging only on
      the path to $\alpha$. TikZ. This is the one figure a viewer of the video
      will remember --- make it good.}
\caption{The GGM-tree construction of a $(2,2)$-distributed point function.}
\label{fig:prim:ggm}
\end{figure}

\subsection{FSS as a comparison gate}
\label{sec:prim:fss-compare}

\todo{Distributed comparison functions and zero-tests evaluated on shared
      inputs in a SINGLE round (\nudgeplain{} Table 3). Grotto (CCS'23) for
      fast comparison over $\Ztwo{n}$. This is what \S4 needs.}

\subsection{Private information retrieval}
\label{sec:prim:pir}

\todo{Multi-server (information-theoretic, needs non-collusion) versus
      single-server (computational, needs LWE and a lot of compute).
      Hafiz--Henry (PoPETs'19.4) is the scheme \pirsonaplain{} builds on, and its
      computationally 1-private variant is DPF-compressed --- so this is the
      SAME primitive as \Cref{sec:prim:fss}, in a different role. Say so
      explicitly; it is the section's thesis.
      Then SimplePIR/DoublePIR (USENIX Sec'23) and Spiral (S\&P'22) for the
      single-server line, and Corrigan-Gibbs--Kogan (EUROCRYPT'20) for the
      preprocessing paradigm that \S5 depends on.}

\subsection{Oblivious RAM}
\label{sec:prim:oram}

\todo{Keep this SHORT --- it matters only for the indexed approaches in \S5.
      Path ORAM (CCS'13), then Floram (CCS'17), Duoram (USENIX Sec'23) and
      PRAC (PoPETs'24). One sentence on why an index and obliviousness are in
      tension: an index is useful precisely because traversal is
      data-dependent, and data-dependent traversal is what obliviousness
      forbids.}

\subsection{Fixed-point arithmetic over rings}
\label{sec:prim:fixedpoint}

\todo{Why this deserves its own subsection rather than a footnote: it is
      where the real engineering cost sits. Represent $v \in \R$ as
      $\lfloor v \cdot 2^{\fracbits} \rfloor \in \ring$; addition is free,
      multiplication needs a TRUNCATION afterwards, and repeated multiplication
      needs NORMALIZATION to stop values growing. Those two operations are the
      entire interactive cost of \S4. Forward-reference it.}
